% ============================================================================
% Part 1 --- Scope, supported profile, reading guide
% ============================================================================
\chapter{Scope and Supported Profile}
\label{ch:scope}

\section{What this document specifies}
\label{sec:what}

This document specifies the \emph{Solidity verifier code generator} located in
\texttt{proofs/\allowbreak{}solidity-verifier}, and justifies that the Solidity/\allowbreak{}Yul program it
emits is a \textbf{faithful translation} of the Rust reference verifier in
\texttt{proofs/\allowbreak{}src/\allowbreak{}plonk/\allowbreak{}verifier.rs} (together with its KZG multi-open
sub-protocol in \texttt{proofs/\allowbreak{}src/\allowbreak{}poly/\allowbreak{}kzg/\allowbreak{}mod.rs}).

``Faithful'' is given a precise operational meaning in
Section~\ref{sec:faithful-meaning}: every algebraic check, every transcript
absorption, every range check, and every rejection path present in the Rust
verifier has a corresponding, semantically equivalent operation in the
generated Solidity, and the two implementations produce \emph{identical}
intermediate values at every semantic checkpoint (verified by differential
tracing, Chapter~\ref{ch:faithfulness}).

\subsection*{Three layers}
The specification tracks three layers throughout:
\begin{enumerate}[leftmargin=1.4em,itemsep=2pt]
  \item \textbf{Rust reference verifier} (\rust{proofs/\allowbreak{}src/\allowbreak{}\ldots}) --- the
        source of truth: halo2-style plonk verification over BLS12-381 with a
        KZG polynomial commitment scheme and a Keccak Fiat--Shamir transcript.
  \item \textbf{Code generator} (\cg{proofs/\allowbreak{}solidity-verifier/\allowbreak{}src/\allowbreak{}\ldots}) ---
        a Rust program that \emph{plans} every fact about the verifier
        (memory layout, proof byte layout, transcript schedule, quotient
        bytecode, PCS schedule) and then renders Solidity/\allowbreak{}Yul through Askama
        templates.
  \item \textbf{Generated Solidity} (\sol{Halo2Verifier.sol}) --- the
        terminal EVM program: a single \texttt{verifyProof(bytes,bytes[])}
        entry point that either returns \texttt{true} or reverts.
\end{enumerate}

\section{The supported profile}
\label{sec:profile}

The generator deliberately supports a \emph{narrow, pinned profile}. Faithfulness
is claimed only inside this profile; outside it the generator refuses to emit
(compile-time or plan-time errors), rather than emitting a subtly wrong verifier.
The profile constraints, and where each is enforced, are:

\begin{longtable}{@{}p{0.30\linewidth}p{0.34\linewidth}p{0.28\linewidth}@{}}
\toprule
\textbf{Constraint} & \textbf{Enforcement site} & \textbf{Rust counterpart} \\
\midrule
\endfirsthead
\toprule
\textbf{Constraint} & \textbf{Enforcement site} & \textbf{Rust counterpart} \\
\midrule
\endhead
Exactly one \emph{committed} instance column (the identity commitment) and
exactly one non-committed instance column.
&
\cg{builder/\allowbreak{}api.rs::validate\_\allowbreak{}instance\_\allowbreak{}column\_\allowbreak{}shape} rejects any other shape.
&
\rust{verifier.rs::verify\_\allowbreak{}algebraic\_\allowbreak{}constraints} instance validation
(committed non-empty, equal lengths, committed $+$ normal $=$
\texttt{num\_\allowbreak{}instance\_\allowbreak{}columns}). \\
\addlinespace
No rotated queries against non-committed instance columns.
&
\cg{builder/\allowbreak{}api.rs} rejects \texttt{RotatedInstanceQuery}.
&
Rust reads non-committed instances only via Lagrange interpolation at
rotation $0$; a rotated read would need a commitment the profile does not have. \\
\addlinespace
BLS12-381 scalar field $\Fr$ with EIP-2537 precompiles; Cancun+ (MCOPY).
&
\cg{evm.rs} pinned \texttt{solc 0.8.30}; \cg{PrecompileSmoke.sol} probes
MCOPY/\allowbreak{}G1ADD/\allowbreak{}G1MSM/\allowbreak{}PAIRING at construction.
&
\rust{transcript/\allowbreak{}implementors.rs} Keccak + EIP-2537 uncompressed G1 encoding. \\
\addlinespace
Feature flags: \texttt{truncated\allowbreak{}challenges},
\texttt{fewer\allowbreak{}point\allowbreak{}sets},
\texttt{outer\allowbreak{}single\allowbreak{}h\allowbreak{}commitment},
\texttt{committed\allowbreak{}instances},
\texttt{keccak\allowbreak{}transcript}.
&
\cg{lib.rs} re-exports; lowering reads them to pick schedules.
&
\rust{poly/\allowbreak{}kzg/\allowbreak{}mod.rs} \texttt{\#[cfg(feature = \ldots)]} branches. \\
\addlinespace
Pinned verifying key: VK digest, payload length, and payload codehash are
baked into the generated contract as constants.
&
\cg{lowering/\allowbreak{}vk.rs::generate\_\allowbreak{}base\_\allowbreak{}vk}; constants
\texttt{EXPECTED\_\allowbreak{}VK\_\allowbreak{}CODEHASH} etc.\ in \sol{Constants.sol}.
&
\rust{plonk/\allowbreak{}mod.rs::hash\_\allowbreak{}into} absorbs \texttt{transcript\_\allowbreak{}repr} first. \\
\bottomrule
\end{longtable}

\begin{remark}[Why a narrow profile is a feature, not a limitation]
Faithfulness of a general lowering of an arbitrary halo2 verifier to gas-costed
Yul is an open-ended problem. By pinning the circuit family (one identity
commitment column, one public-instance column, fixed feature flags, pinned VK),
the space of generated programs becomes finite and fully inspectable: every
generated contract in this repository is accompanied by a recorded
\texttt{EXPECTED\_\allowbreak{}VK\_\allowbreak{}CODEHASH} and a differential trace test that pins every
intermediate scalar. The audit surface collapses from ``all circuits'' to
``the lowering function itself''.
\end{remark}

\section{Reading guide}
\label{sec:reading}

\begin{itemize}[leftmargin=1.4em,itemsep=2pt]
  \item \textbf{Chapter~\ref{ch:math}} states the verifier algorithm as
        mathematics: transcript sampling, Lagrange evaluation, the four
        identity families (gate, permutation, lookup, trash), the
        $y$-batched quotient numerator, the linearization, the KZG multi-open
        reduction, and the final pairing. Every later section refers back to
        equation numbers here.
  \item \textbf{Chapter~\ref{ch:rust}} walks the Rust reference verifier
        function by function, with exact code excerpts.
  \item \textbf{Chapter~\ref{ch:codegen}} walks the code generator: how it
        plans the protocol, the memory map, the proof layout, the quotient
        virtual machine, and the PCS schedule.
  \item \textbf{Chapter~\ref{ch:solidity}} walks the generated Solidity/\allowbreak{}Yul.
  \item \textbf{Chapter~\ref{ch:correspondence}} is the core deliverable:
        per-stage tables mapping \rust{Rust} $\leftrightarrow$ \cg{codegen}
        $\leftrightarrow$ \sol{Solidity} with equation references.
  \item \textbf{Chapter~\ref{ch:faithfulness}} assembles the faithfulness
        argument: semantic checkpoints, trace-equivalence testing, mutation
        testing, and an explicit list of every check in the Rust verifier and
        where it lives in the generated code.
  \item \textbf{Chapter~\ref{ch:ledger}} is the coverage ledger: every module
        and public function visited, with status.
\end{itemize}

\subsection*{Notation conventions}
Throughout, \rust{orange text} denotes Rust reference-verifier identifiers,
\cg{blue text} denotes code-generator identifiers, and \sol{slate text}
denotes generated Solidity/\allowbreak{}Yul identifiers. Trace identifiers are written
\boxid{id~N}. Memory pointers are written in hexadecimal, e.g.\
\texttt{0x2b60}.
